\documentclass[11pt]{article}

\usepackage[margin=1in]{geometry}
\usepackage{amsmath,amssymb,booktabs,tabularx}

\title{Derivation title}
\author{}
\date{}

\begin{document}
\maketitle

\noindent\textbf{Status.} Draft\\
\textbf{Plan item.} Exact derivation name from \texttt{design/paper/plan.md}\\
\textbf{Depends on.} None or linked derivation files\\
\textbf{Feeds.} Paper section, appendix, proposition, plot, or claim

\section{Purpose}

State the specific paper need this derivation satisfies. Keep this focused on
the mathematical role, not the implementation.

\section{Setup}

Define the time interval, conditioning information, state variables, parameters,
and random objects used in the derivation.

Use this section to make all conventions explicit:
\begin{itemize}
    \item Conditioning sigma-field or initial state.
    \item Measure used in the derivation.
    \item Clock, seasonality, or schedule convention.
    \item Normalization convention.
    \item Domain restrictions needed before any transform is evaluated.
\end{itemize}

\section{Notation}

\begin{tabularx}{\textwidth}{lXX}
\toprule
Symbol & Meaning & Units or domain\\
\midrule
\(\cdots\) & \(\cdots\) & \(\cdots\)\\
\bottomrule
\end{tabularx}

\section{Assumptions used}

\begin{tabularx}{\textwidth}{XXX}
\toprule
Assumption & Where it is used & Why it is needed\\
\midrule
\(\cdots\) & \(\cdots\) & \(\cdots\)\\
\bottomrule
\end{tabularx}

\section{Result}

State the final result in the form that should appear in the paper.

If the result is a proposition, use:

\medskip
\noindent\textbf{Proposition.} \(\cdots\)

\medskip
If the result is an identity or construction, use:

\medskip
\noindent\textbf{Identity.} \(\cdots\)

\section{Derivation}

Work through the derivation step by step.

Keep each algebraic transition justified. If a step depends on independence,
conditioning, affine closure, a moment-domain constraint, or a normalization,
name that reason explicitly.

\section{Domain and validity}

List the exact conditions under which the result is valid.

Include transform domains, positivity constraints, stability requirements,
integrability requirements, and any terminal or boundary conditions.

\section{Consequences}

Record the direct implications that should be reused elsewhere in the paper.

\begin{itemize}
    \item Consequence 1.
    \item Consequence 2.
\end{itemize}

\section{Checks}

Describe how the result should be checked analytically or numerically.

\begin{tabularx}{\textwidth}{XXX}
\toprule
Check & Expected result & Evidence location\\
\midrule
\(\cdots\) & \(\cdots\) & \(\cdots\)\\
\bottomrule
\end{tabularx}

\section{Open questions, if any}

\begin{itemize}
    \item \(\cdots\)
\end{itemize}

\end{document}
